\documentclass[twocolumn]{article}

% Essential packages
\usepackage{hyperref}
\usepackage{graphicx}
\usepackage{amsmath}
\usepackage{amssymb}
\usepackage{booktabs}
\usepackage[authoryear,round]{natbib}
\usepackage{appendix}


% Page geometry
\usepackage[letterpaper,
            includeheadfoot,
            layoutsize={8.125in,10.875in},
            layouthoffset=0.1875in,
            layoutvoffset=0.0625in,
            left=38.5pt,
            right=43pt,
            top=43pt,
            bottom=44pt,
            headheight=0pt,
            headsep=-2pt,
            columnsep=0.3125in,
            footskip=25pt,
            marginparwidth=38pt]{geometry}

% Authors and affiliations
\usepackage{authblk}

% Float control
\renewcommand{\topfraction}{0.9}
\renewcommand{\bottomfraction}{0.9}
\renewcommand{\textfraction}{0.05}
\renewcommand{\floatpagefraction}{0.9}

% Math operators
\DeclareMathOperator*{\mean}{mean}
\DeclareMathOperator*{\std}{std}
\DeclareMathOperator*{\AbsDiff}{AbsDiff}

% Bibliography settings
\setlength{\bibsep}{0pt plus 0.3ex}

% First page footer
\usepackage{fancyhdr}
\fancypagestyle{firstpage}{%
  \fancyhf{}%
  \fancyfoot[C]{\footnotesize
    This version has not been peer reviewed. Comments and feedback are welcome.}%
  \renewcommand{\headrulewidth}{0pt}%
}

% Caption formatting
\usepackage{caption}
\captionsetup{labelfont=bf, font=small}

% Author/affiliation fonts
\renewcommand\Authfont{\large}
\renewcommand\Affilfont{\normalsize}

% Autoref customization
\renewcommand{\tableautorefname}{Tab.}
\renewcommand{\figureautorefname}{Fig.}
\renewcommand{\equationautorefname}{Eq.}
\renewcommand{\appendixautorefname}{Appx.}

% TODO: confirm Jorge's affiliation number and institution
\title{\LARGE Probing Warmth and Competence Representations in LLM Hiring Decisions}
\date{\today}
\author[1]{Emrecan Ulu}
\author[2]{Jorge Lastra Cerda}
\affil[1]{University of Konstanz}
\affil[2]{University of Konstanz}  % TODO: update if different

\begin{document}
\twocolumn[
  \begin{@twocolumnfalse}
\maketitle
\thispagestyle{firstpage}


\begin{abstract}
\noindent
% TODO: write abstract once results are finalized.
% Suggested structure (3--4 sentences):
% (1) Discrimination in hiring is well-documented in humans and increasingly in AI.
% (2) We probe whether warmth and competence (Stereotype Content Model) are
%     encoded as linear directions in the residual stream of Gemma 3 12B.
% (3) We steer those directions and measure causal shifts in callback recommendations,
%     benchmarked against Gallo and Hausladen (2024).
% (4) Findings and implications.
\texttt{Due: Week 8}
\vspace{0.50cm}
\end{abstract}
  \end{@twocolumnfalse}
]


\renewcommand{\thefootnote}{\fnsymbol{footnote}}
\footnotetext[0]{Corresponding author: TODO@uni-konstanz.de}  % TODO: update email
\renewcommand{\thefootnote}{\arabic{footnote}}
\setcounter{footnote}{0}

% ─────────────────────────────────────────────────────────────────────────────
\section*{Introduction}
% ─────────────────────────────────────────────────────────────────────────────

Discrimination in hiring is one of the most persistent and consequential
forms of social inequality.
Access to work is not a neutral outcome: it shapes social mobility,
material security, and the possibilities available to people across their
lives.
One of the most cited field experiments documenting the scope of this
problem is \citet{bertrand2004emily}, who sent nearly 5{,}000 matched
resumes to employers in Boston and Chicago, varying only the perceived
race of the applicant through the name on the resume.
Resumes assigned White-sounding names received approximately 50~percent
more callbacks than otherwise identical resumes assigned
African-American-sounding names, a gap that held across industries,
occupations, and firm sizes.
Similar audits have since documented comparable patterns for gender, age,
disability, religion, sexual orientation, and national origin,
establishing that differential treatment in hiring is not an isolated
phenomenon but a structural feature of labor markets in North America and
beyond \citep{correll2007getting, neumark2019older, tilcsik2011pride}.

Social psychology offers a theoretical framework for understanding why
these disparities exist and what drives their magnitude.
The Stereotype Content Model \citep[SCM;][]{fiske2002model} proposes that
social perception is organized along two fundamental dimensions: warmth,
which captures perceived intentions and trustworthiness, and competence,
which captures perceived ability and effectiveness.
Groups stereotyped as low on both dimensions elicit contempt and active
harm, while groups stereotyped as high on both elicit admiration and
active facilitation \citep{cuddy2007biasmap}.
\citet{gallo2024warmth} tested this model directly against hiring outcomes.
In a meta-analysis pooling callback data from over two decades of North
American correspondence studies, they showed that human warmth and
competence ratings of social signals predict callback disparities with
substantial accuracy.
The way people perceive the warmth and competence of a social group is
not just a psychological description: it is a working predictor of
whether members of that group get called back for a job interview.

Large language models are now entering this picture.
They are being used to screen resumes, rank candidates, and support or
automate hiring decisions at scale.
This raises an urgent question: do these systems reproduce the
discriminatory patterns documented in human hiring, and if so, what is
the underlying mechanism?
Behavioral audit studies have already begun to answer the first part.
\citet{an2024llm} found systematic disparities in GPT-3.5 callback
recommendations across racial and ethnic groups.
\citet{gaebler2024auditing} audited multiple large language models and found
persistent race and gender disparities with potential legal implications
under equal employment opportunity law.
What these studies cannot tell us is where in the model this differential
treatment originates, and how it could be addressed at the source rather
than patched at the output.

% ─────────────────────────────────────────────────────────────────────────────
\section*{Literature}
% ─────────────────────────────────────────────────────────────────────────────

Four bodies of literature converge on the problem this paper addresses,
and each leaves a specific gap.

\subsection*{Correspondence Studies}

Correspondence studies have proven that discrimination in hiring is real
and measurable across race, gender, age, disability, and national origin.
What they cannot tell us is why it happens at the level of cognitive
mechanism: they describe the outcome but not the process that produces it.

\subsection*{The Stereotype Content Model}

The SCM answers that question for humans.
Warmth and competence are the two dimensions through which social groups
are perceived, and those perceptions predict who gets called back, as
\citet{gallo2024warmth} demonstrated directly.
But the SCM was built entirely within human psychology.
Nobody has tested whether this same structure exists inside a machine.
If a language model trained on human-generated text has internalized
warmth and competence as internal representations, that would be a
meaningful finding; if it has not, that is equally important to know.

\subsection*{Behavioral Audits of Large Language Models}

Behavioral audits of large language models have confirmed that AI systems
also discriminate in hiring \citep{an2024llm, gaebler2024auditing}.
But by treating the model as a black box, these studies can document that
disparities exist without explaining where in the model they originate or
what internal structure produces them.
Knowing that a model recommends against an applicant with a Black-sounding
name more often than an identical applicant with a White-sounding name is
useful, but it does not tell us how to fix it or where the bias lives.

\subsection*{Mechanistic Interpretability}

Mechanistic interpretability gives us tools to open that box.
Representation engineering \citep{zou2023repeng} and activation addition
\citep{turner2023activation} have shown that high-level concepts are encoded as
linear directions in large language model residual streams and can be
causally manipulated.
Most recently, \citet{sofroniew2026emotion} applied this framework to emotion
concepts in Claude Sonnet, showing that internal emotion representations
causally influence model behavior.
But these tools have never been applied to warmth and competence as social
stereotype dimensions, and never connected to hiring discrimination.

\subsection*{This Paper}

This paper sits at the intersection of all four gaps.
We take the conceptual framework from the SCM, the empirical benchmark
from correspondence studies, the motivating question from LLM bias audits,
and the methodological tools from mechanistic interpretability.
We open Gemma~3 12B and look for internal vectors of warmth and
competence, to test whether the SCM structure is not just a human
phenomenon but something the model has learned from human-generated text,
and whether those representations causally drive the hiring disparities
that behavioral audits have documented from the outside.

% ─────────────────────────────────────────────────────────────────────────────
\section*{Methods}
% ─────────────────────────────────────────────────────────────────────────────
\texttt{Due: Week 10.}

\paragraph{Story Corpus.}
% TODO: describe the 4-condition story generation (high/low warmth x high/low competence,
% 100 topics x 12 stories = 4,800 total) using claude-haiku-4-5 via the Anthropic API.

\paragraph{Activation Extraction.}
% TODO: describe TransformerLens hook at blocks.31.hook\_resid\_post,
% token averaging from position 50 onward, contrast vector construction.

\paragraph{PCA Denoising.}
% TODO: describe neutral corpus, PCA, projection-out of top components
% explaining 50\% of variance.

\paragraph{Probe Validation.}
% TODO: describe cross-validated linear probe and logit-lens check.

\paragraph{Causal Steering.}
% TODO: describe steering strengths (0, 0.25, 0.5, 1.0 x mean residual norm),
% hiring prompt template, Yes/No logit scoring.

\paragraph{Benchmark.}
We compare model callback disparities against the warmth and competence
ratings and callback rates from \citet{gallo2024warmth}, which provides
empirical grounding in human correspondence studies.

\paragraph{Data.}
Hiring stimuli are built around the 282 applicant names and category-level
signals from \citet{gallo2024warmth}, paired with a fixed administrative
assistant resume template to isolate the effect of the social signal.

% ─────────────────────────────────────────────────────────────────────────────
\section*{Results}
% ─────────────────────────────────────────────────────────────────────────────
\texttt{Due: Week 10.}

% Suggested structure for 3 main results:
%
% \paragraph{Warmth and Competence Are Linearly Encoded in Gemma 3 12B.}
% Report probe accuracy, Cohen's d, logit-lens semantic check.
%
% \paragraph{Steering Warmth and Competence Causally Shifts Callback Recommendations.}
% Report logit delta as a function of steering strength.
%
% \paragraph{Model Callback Disparities Mirror the Human Benchmark.}
% Report correlation between model disparity and Gallo--Hausladen callback rates.

% Preliminary smoke-test result (early technical check, move to results when full pipeline runs):
A cross-validated linear probe on warmth versus cold sentences achieves
86~percent accuracy at layer~31 of Gemma~3 12B (Cohen's $d = 2.90$).
Warmth steering produces roughly twice the logit shift of an equivalent
random direction, consistent with the vector encoding something beyond
surface-level noise.

% ─────────────────────────────────────────────────────────────────────────────
\section*{Discussion and Conclusion}
% ─────────────────────────────────────────────────────────────────────────────
\texttt{Due: Week 11.}

If the warmth and competence vectors we extract correspond to the same
dimensions that drive discrimination in human hiring, this has implications
beyond documenting another case of algorithmic bias.
It suggests a mechanism: models may have learned an internal structure of
social perception that mirrors the SCM, inherited from the human-generated
text on which they were trained.
It opens a path to debiasing at the source, targeting the representation
rather than filtering outputs after the fact.
And it provides a basis for understanding how AI hiring tools can fail and
potentially be exploited.

\paragraph{Limitations.}
% TODO: smoke-test stimuli have warmth-valence confound (addressed by topic control
% and PCA denoising in the main pipeline but worth noting);
% results are model-specific (Gemma 3 12B); hiring prompt is simplified.

\paragraph{Future Work.}
% TODO: scale to Gemma 3 27B; test across additional models (Llama 3, Qwen 2.5);
% extend to category-level signals beyond names; connect to SAE feature decomposition.

{
\small

\clearpage
\bibliographystyle{apalike}  % CHANGED: matches authoryear natbib mode
\bibliography{references}
}



\clearpage
\begin{appendices}
\onecolumn
\setcounter{figure}{0}
\setcounter{table}{0}
\renewcommand{\thesection}{\Alph{section}}%
\renewcommand{\thesubsection}{\thesection.\arabic{subsection}}
\renewcommand\thefigure{S.\arabic{figure}}
\renewcommand\thetable{S.\arabic{table}}
\begin{center}
    \begin{minipage}{\textwidth}
        \centering
        \fontsize{16}{20}\selectfont
        Supplementary Materials
    \end{minipage}
\end{center}
\vspace{.5cm}

\section*{Details on Methods}
\label{apx:methods}

\section*{Additional Results}

\end{appendices}
\end{document}
